\centerline{\bf \Huge Пробный ЕГЭ II}
\centerline{\bf \Huge }

\begin{flushright}
        \scriptsize{\textbf{94+}}
\end{flushright}

\problem{1} Четырёхугольник $ABCD$ вписан в окружность. Угол $ABC$ равен 107°, угол $CAD$ равен 43°. Найдите угол $ABD$. Ответ дайте в градусах.

\problem{2} Даны векторы $\vec{a}(3;0)$ и $\vec{b}(1;3)$. Найдите длину вектора $\vec{a} +5 \vec{b}$.

\problem{3} Найдите объём многогранника, вершинами которого являются вершины $A, B, C, D, B_1$ прямоугольного параллелепипеда $ABCDA_1B_1C_1D_1$, у которого $AB=10, BC=4, BB_1=9$.

\problem{4} В группе туристов 20 человек. С помощью жребия они выбирают семь человек, которые должны идти в село в магазин за продуктами. Какова вероятность того, что турист Д, входящий в состав группы, пойдет в магазин?

\problem{5} Стрелок стреляет по одному разу в каждую из четырёх мишеней. Вероятность попадания в мишень при каждом отдельном выстреле равна 0,9. Найдите вероятность того, что стрелок попадёт в первую мишень и не попадёт в три последние.

\problem{6} Найдите корень уравнения $\sqrt{53-2x}=7$.


\problem{7} Найдите значение выражения $3\sin \dfrac{13 \pi}{12} \cdot \cos \dfrac{13 \pi}{12}$.

\problem{8} На рисунке изображен график $y=f^{\prime}(x)$ - производной функции $f(x)$, определенной на интервале ( $-15 ; 5$ ). Найдите количество точек максимума функции $f(x)$, принадлежащих отрезку $[-11 ;4]$.

\problem{9} Мотоциклист, движущийся по городу со скоростью $v_0=95 \mathrm{км} / \mathrm{ч}$, выезжает из него и сразу после выезда начинает разгоняться с постоянным ускорением $a=40 \mathrm{км} / \mathrm{ч}^2$. Расстояние (в км) от мотоциклиста до города вычисляется по формуле $S=v_0 t-\dfrac{a t^2}{2}$, где $t$ - время в часах, прошедшее после выезда из города. Определите время, прошедшее после выезда мотоциклиста из города, если известно, что за это время он удалился от города на 25 км. Ответ дайте в минутах.


\problem{10} Поля и Оля, работая вместе, пропалывают грядку за 18 минут, а одна Оля - за 30 минут. За сколько минут пропалывает эту грядку одна Поля?


\problem{11} На рисунке изображены графики функций $f(x)=a^x$. Найдите значение $f(4)$.

\problem{12} Найдите точку максимума функции  $y=8\cdot ln(x-5) -8x+9$. 

\problem{13} a) Решите уравнение $\cos 2x - \sin(x-\pi)-1=0$.

б) Укажите корни этого уравнения, принадлежащие отрезку $\left[-\dfrac{7 \pi}{2} ;-2\pi\right]$.

\problem{14} В правильной четырёхугольной пирамиде $S A B C D$ с основанием $ABCD$ точка $O$ - центр основания пирамиды, точка $M$ - середина ребра $SC$, точка $K$ делит ребро $BC$ в отношении $BK:KC=3:2$, а $AB=4$ и $SO=2\sqrt{23}$.

a) Докажите, что плоскость $OMK$ параллельна прямой $SA$.

б) Найдите длину отрезка, по которому плоскость $OMK$ пересекает грань $SAD$.

\problem{15} Решите неравенство $\dfrac{2^x+8}{2^x-8}+\dfrac{2^x-8}{2^x+8} \geqslant \dfrac{2^{x+4}+96}{4^x-64}$.


\problem{16} В июле 2026 года планируется взять кредит на некоторую сумму. 

Условия возврата таковы:

- каждый январь долг увеличивается на $20 \%$ по сравнению с концом предыдущего года;

- с февраля по июнь каждого года необходимо выплатить одним платежом часть долга.

Сколько рублей планируется взять в банке, если известно, что кредит будет полностью погашен тремя равными платежами (то есть за три года) и общая сумма выплат после полного погашения кредита на 77200 рублей больше суммы, взятой в кредит?

\problem{17} Пятиугольник $A B C D E$ вписан в окружность. Известно, что $A B=C D=3$ и $B C=D E=4$.

a) Докажите, что $A C=C E$.

б) Найдите длину диагонали $B E$, если $A D=6$.

\problem{18} Найдите все значения параметра $a$, при каждом из которых система уравнений

$$
\left\{\begin{array}{l}
y=|x-a|-1 \\
|y|+x^2-2x=0
\end{array}\right.
$$


имеет ровно четыре различных решения.

\problem{19} В порту имеются только заполненные контейнеры, масса каждого из которых равна 20 тонн или 40 тонн. В некоторых контейнерах находится сахарный песок. Количество контейнеров с сахарным песком составляет $60 \%$ от общего числа контейнеров.

a) Может ли масса контейнеров с сахарным песком составлять $50 \%$ от общей массы?

б) Может ли масса контейнеров с сахарным песком составлять $40 \%$ от общей массы?

в) Какую наибольшую долю (в процентах) может составлять масса контейнеров с сахарным песком от общей массы всех контейнеров?

